% ==========================================================================
%  Thesis Notes — Kafka + R: Findings, Limitations, and Future Work
%  For inclusion/reference in the final thesis document
% ==========================================================================
\documentclass[11pt,a4paper]{article}
\usepackage[margin=2.5cm]{geometry}
\usepackage{booktabs,tabularx}
\usepackage{hyperref}
\usepackage[T1]{fontenc}
\usepackage{lmodern}
\usepackage{listings}
\lstset{basicstyle=\ttfamily\small, breaklines=true}

\title{Thesis Notes\\[4pt]\large REST vs Kafka Architecture Comparison in R Shiny}
\author{Iva Gunjaca — University of Amsterdam}
\date{June 2026}

\begin{document}
\maketitle

% ------------------------------------------------------------------
\section{Core Finding: R's Runtime is the Bottleneck}
% ------------------------------------------------------------------

R Shiny runs all application logic---UI rendering, reactive computations,
and Kafka message consumption---on a single thread. The \texttt{reactivePoll()}
mechanism fires a callback at a fixed interval (e.g.\ every 150\,ms), and
\texttt{consumer\$consume(timeout)} blocks that thread for up to
\texttt{timeout} milliseconds waiting for a message. During this blocking
window, no UI updates, plot rendering, or other reactive computations can
execute.

This creates a fundamental conflict: lowering the consume timeout improves
UI responsiveness but reduces the chance of catching a message per cycle;
raising it improves message throughput but starves the UI. No configuration
of these two parameters resolves the underlying issue.

% ------------------------------------------------------------------
\section{The R Kafka Ecosystem}
% ------------------------------------------------------------------

As of June 2026, only two R packages for Apache Kafka exist:

\begin{tabularx}{\textwidth}{@{}llX@{}}
\toprule
\textbf{Package} & \textbf{Status} & \textbf{Mechanism} \\
\midrule
\texttt{kafka} (INWTlab/r-kafka) & Active, not on CRAN &
  C wrapper around librdkafka via Rcpp. Exposes synchronous
  \texttt{consume(timeout\_ms)} and \texttt{produce(topic, value, key)}.
  No async API, no callbacks, no batch consumption. \\
\texttt{rkafka} & Archived from CRAN (May 2023) &
  rJava wrapper requiring a JVM inside the R process.
  Removed for unresolved build issues. Deprecated. \\
\bottomrule
\end{tabularx}

Neither package provides the primitives needed for real-time streaming:
asynchronous message delivery, background thread consumption, or
event-driven callbacks. This is not an oversight---R is not used for
streaming workloads, so the demand for such libraries does not exist.

% ------------------------------------------------------------------
\section{Why More Users Increases Latency}
% ------------------------------------------------------------------

All Kafka consumers in the Shiny frontends subscribe to a shared
\texttt{output} topic. When multiple users are active:

\begin{enumerate}
  \item Each user's consumer reads from the same topic partition(s).
  \item Each \texttt{consume()} call may return a message intended for a
        different session or app, which is discarded.
  \item The consumer must poll again (waiting another full cycle) to find
        its own message.
  \item Meanwhile, the R backend processes \texttt{input} messages
        sequentially, falling further behind under concurrent load.
\end{enumerate}

This is a fan-out problem on a shared topic. With $N$ concurrent guests,
the probability of receiving an irrelevant message per poll increases,
compounding latency.

In the REST architecture, this problem does not exist: each GET request
performs an O(1) Redis lookup by session ID. No scanning, no shared queue.

% ------------------------------------------------------------------
\section{Why Separate Containers Do Not Help}
% ------------------------------------------------------------------

Shiny hosting platforms (shinyapps.io, Shiny Server Pro) isolate users
into separate R processes. This helps when the bottleneck is CPU contention
between users. However, the Kafka consumer bottleneck is per-session,
not per-pod:

\begin{itemize}
  \item Even with a single user in a dedicated container, the
        \texttt{consume()} call still blocks the R event loop.
  \item Heavyweight apps (MonteCarlo, ClimateAnomaly, MLTrainer) leave
        less time for the consumer between UI updates due to heavier
        R processing (simulation loops, model training, CSV parsing).
\end{itemize}

% ------------------------------------------------------------------
\section{Why a Better R Kafka Library Cannot Fully Solve This}
% ------------------------------------------------------------------

Even a hypothetical async Kafka library for R would face R's fundamental
limitation: the interpreter is single-threaded. An async library could
buffer messages on a background thread, but the R code that processes
those messages (deserializing JSON, updating reactive values, triggering
UI re-renders) still runs on the main thread and cannot be parallelized.

The \texttt{httpuv} package (which powers Shiny's HTTP server) demonstrates
this pattern: it runs a C++ event loop on a background thread for I/O,
but all R callback execution is serialized on the main thread. A Kafka
library could adopt the same design, but the throughput ceiling would
remain R's single-threaded processing speed.

% ------------------------------------------------------------------
\section{Latency vs Data Integrity: The Case for Kafka Despite Higher Latency}
% ------------------------------------------------------------------

\subsection{What Kafka Provides That REST Does Not}

\begin{tabularx}{\textwidth}{@{}lXX@{}}
\toprule
\textbf{Property} & \textbf{REST + Redis} & \textbf{Kafka} \\
\midrule
State durability
  & Last-write-wins; intermediate states are overwritten and lost
  & Append-only log; every state change is persisted \\
Ordering guarantees
  & None; concurrent writes may arrive out of order
  & Guaranteed per partition \\
Crash recovery
  & State lost if Redis restarts without persistence
  & Consumer resumes from last committed offset \\
Audit trail
  & Must be explicitly implemented (e.g.\ event-sourcing table)
  & Built-in via topic retention policy \\
Data loss under concurrency
  & Possible: if two users write between polls, one update is silently
    overwritten
  & Impossible: all messages are appended, none overwritten \\
\bottomrule
\end{tabularx}

\subsection{When Latency Does Not Matter}

For scientific use cases where the workflow is:

\begin{enumerate}
  \item Researcher A configures simulation parameters,
  \item A long-running computation executes (minutes to hours),
  \item Researcher B reviews the results,
\end{enumerate}

\noindent the relay latency is irrelevant. What matters is that the
parameter configuration is \textbf{not lost or corrupted} between steps.
Kafka's append-only log guarantees this: every parameter change is
recorded in order, and if the system crashes mid-computation, the exact
state can be reconstructed from the topic.

In contrast, the REST architecture stores only the \textit{latest} state
in Redis. If Redis restarts, or if two researchers modify parameters
simultaneously, intermediate states are silently lost. For a
reproducibility-focused scientific workflow, this is a more serious
deficiency than high latency.

\subsection{When Latency Does Matter}

For interactive collaborative editing---multiple researchers adjusting
a shared visualization in real time---latency directly impacts usability.
At multi-second delays, users make conflicting changes based on stale
state, may assume the platform is broken, and retry actions, creating
duplicate state changes.

\subsection{Synthesis}

The choice between REST and Kafka is not ``which is better'' but
``which trade-off is acceptable for the use case'':

\begin{itemize}
  \item \textbf{Interactive collaboration} (real-time parameter tuning,
        shared dashboards): REST wins. Low latency is essential, and
        data integrity can be supplemented with explicit save/restore
        snapshots to PostgreSQL.
  \item \textbf{Asynchronous scientific workflows} (long-running
        simulations, reproducible pipelines): Kafka wins. Latency is
        irrelevant, and the append-only log provides an audit trail
        and crash recovery that REST cannot match without additional
        infrastructure.
  \item \textbf{Hybrid approach}: use REST for the real-time state
        synchronization layer (low latency for UI updates) and Kafka
        as a durable event log for audit, replay, and crash recovery.
        This combines the strengths of both architectures but increases
        operational complexity.
\end{itemize}

% ------------------------------------------------------------------
\section{Benchmark Load Justification}
% ------------------------------------------------------------------

The k6 benchmarks use the following virtual user (VU) counts:

\begin{tabularx}{\textwidth}{@{}lrX@{}}
\toprule
\textbf{Benchmark} & \textbf{Peak VUs} & \textbf{Ramp pattern} \\
\midrule
01 State Relay          & 20  & 1 $\to$ 5 $\to$ 10 $\to$ 20 $\to$ 10 $\to$ 1 \\
02 Collaboration        & 30  & 1 $\to$ 10 $\to$ 20 $\to$ 30 $\to$ 15 $\to$ 1 \\
03 Save/Restore         & 15  & Ramped \\
04 Throughput           & 100 & 0 $\to$ 10 $\to$ 25 $\to$ 50 $\to$ 75 $\to$ 100 $\to$ 50 $\to$ 0 \\
05 Data Loss            & 15  & Fixed \\
06 Cross-Contamination  & ---  & Fixed (2 sessions per app) \\
07 Session Lifecycle    & 50  & Ramped \\
08 Multi-User           & 3   & 3 VUs simulating groups of 3, 5, and 10 guests \\
09 WebSocket Presence   & 50  & Ramped \\
\bottomrule
\end{tabularx}

These loads are realistic for a university collaborative research platform.
Typical usage patterns include:

\begin{itemize}
  \item \textbf{5--20 concurrent users}: a research lab group or seminar
        working on a shared analysis session.
  \item \textbf{50--100 concurrent users}: a lecture demonstration, a
        workshop, or a course exercise where an entire class interacts
        with the platform simultaneously.
\end{itemize}

The single-node Hetzner CX32 cluster (8 vCPU, 16\,GB RAM) is
representative of what a university research group would realistically
deploy.

% ------------------------------------------------------------------
\section{Benchmarking Methodology: Baseline-Calibrated Thresholds}
% ------------------------------------------------------------------

The k6 benchmark suite uses \emph{per-app baseline-calibrated thresholds}
rather than uniform SLAs across all applications.

\subsection{Problem with Uniform Thresholds}

A flat threshold (e.g., ``relay latency p95 < 10\,s'' for all apps)
conflates \emph{architecture overhead} with \emph{computation cost}.
A calculator performing \texttt{10 + 20} and an ML trainer fitting a
500-tree Random Forest on 2000 rows have fundamentally different
baseline response times. Applying the same SLA to both measures
whether R can compute faster, not whether the architecture adds
acceptable overhead.

\subsection{Baseline Calibration Process}

The calibration follows a two-phase approach:

\begin{enumerate}
  \item \textbf{Baseline measurement} (benchmark 00): Run each of the
        8 applications with 1 virtual user, no concurrency, 5 repetitions.
        Record per-app median response times for: POST (state relay),
        GET (cache read), relay round-trip (POST + poll until result
        appears), save-to-DB, and restore-from-DB.

  \item \textbf{Threshold derivation}: Set each app's threshold at
        $5\times$ its baseline median. This multiplier accounts for
        degradation under concurrent load while remaining proportional
        to the app's natural computation cost.
\end{enumerate}

\noindent Thresholds are applied per-app using k6's tagged threshold
syntax: \texttt{metric\{app\_name:Calculator\}: ['p(95)<1000']}.

\subsection{Deployment Consideration: Shared Volume with Replicas}

The ClimateAnomaly application requires a CSV file on a shared volume
(\texttt{/app/shared\_data/}). When running multiple backend replicas,
the PersistentVolumeClaim is declared as \texttt{ReadWriteOnce}, meaning
it can only be mounted by pods on the same node. On a single-node cluster
all replicas share the same PVC, but \texttt{kubectl cp} only targets one
pod. If the K8s Service load-balances a request to a different replica
that has not yet written the file to its local cache, the request fails
with a 500 error.

The solution is to copy the test data to all replicas before running
benchmarks:

\begin{lstlisting}[language=bash]
for POD in $(kubectl get pod -l app=shiny-back-csv-advanced -o name); do
  kubectl cp k6/bench-climate-data.csv \
    ${POD#pod/}:/app/shared_data/bench-climate-data.csv
done
\end{lstlisting}

\noindent This is a test environment concern, not an architectural
limitation. In production, the CSV would be uploaded via the application
UI and written to the shared volume by the receiving pod, making it
available to all replicas through the PVC.

\subsection{Bug: Output File Overwrites Input File}

The ClimateAnomaly backend writes its processed summary CSV back to the
shared volume. The output filename was derived using
\texttt{sub("\^{}raw\_", "processed\_", body\$file)}, which only adds the
\texttt{processed\_} prefix if the input filename starts with
\texttt{raw\_}. The benchmark test file \texttt{bench-climate-data.csv}
does not, so the output was written to the same path as the input,
overwriting it.

On the first request, the backend reads the 98-row input CSV, processes
it into a 5-row summary, and writes the summary back to
\texttt{bench-climate-data.csv}. On the second request, the backend reads
the now-5-row summary file, attempts \texttt{as.Date(df\$Timestamp)} on
a column that no longer exists, and crashes with
\texttt{replacement has 0 rows, data has 5}.

The fix: always prefix the output filename with \texttt{processed\_}
using \texttt{paste0("processed\_", body\$file)} instead of \texttt{sub()}.
This ensures the input file is never overwritten regardless of its naming
convention.

\subsection{What the Baseline Separates}

The per-app threshold design separates two distinct latency components:

\begin{description}
  \item[Computation latency:] Time spent in the R backend executing
    application logic (e.g., \texttt{randomForest()}, CSV parsing,
    anomaly detection). This is inherent to the workload and cannot
    be reduced by the messaging architecture.

  \item[Architecture latency:] Time added by the state relay pipeline
    (Spring Boot routing, Redis caching, Plumber HTTP overhead, or
    Kafka produce/consume cycle). This is what the thesis compares
    between REST and Kafka.
\end{description}

\noindent By setting thresholds relative to each app's baseline, the
benchmarks measure whether the architecture adds \emph{acceptable
overhead on top of the app's natural cost}, rather than whether R
can compute faster. A test failure under this scheme indicates
architecture-induced degradation, not workload complexity.

\subsection{Queuing-Based Multiplier Justification}

The degradation factor accounts for request serialization at the R
backend. R's Plumber server processes requests sequentially on a single
thread, forming an M/M/1-like queue. With $P$ concurrent users routed
to one replica, the expected worst-case latency is approximately
$P \times \text{baseline}$.

With 3 replicas and varying VU counts per test:

\begin{center}
\begin{tabular}{lrrrl}
\toprule
\textbf{Test} & \textbf{Peak VUs} & \textbf{Per replica} & \textbf{Multiplier} & \textbf{Rationale} \\
\midrule
01, 02, 07, 08 & 20--50 & $\sim$7--17 & 10$\times$ & Moderate queuing \\
04 Throughput   & 100    & $\sim$33    & 50$\times$ & Heavy queuing \\
03 Save/Restore & 15     & $\sim$5     & 10$\times$ & PostgreSQL parallel, pool-limited \\
\bottomrule
\end{tabular}
\end{center}

\noindent The multiplier is intentionally generous---the goal is not to
enforce a production SLA, but to detect \emph{disproportionate}
architectural overhead beyond normal queuing. A test failure at
$10\times$ baseline indicates a systemic issue (connection exhaustion,
memory pressure, thread starvation), not expected load degradation.

For compute-intensive apps like MLTrainer (baseline POST: 14.2\,s),
the resulting threshold (711\,s for throughput) is effectively
``never fail.'' This is correct: the architecture overhead is
negligible compared to the 14-second RandomForest training time.
The test confirms that the architecture does not add disproportionate
overhead to compute-intensive workloads---a valid finding, not a
testing gap.

\subsection{References}

\begin{itemize}
  \item \textbf{Harchol-Balter (2013)} provides the queuing theory
        foundation for understanding request serialization in
        single-threaded servers. Chapter~13 (M/M/1 and PASTA) directly
        models the behavior of R's Plumber backend: a single server
        processing requests sequentially with Poisson arrivals. The
        expected response time grows linearly with utilization, which
        justifies the per-VU multiplier approach.\\
        \emph{Reference:} Harchol-Balter, M., \emph{Performance
        Modeling and Design of Computer Systems: Queueing Theory in
        Action}, Cambridge University Press, 2013.
        ISBN: 978-1-107-02750-3.

  \item \textbf{Jain (1991)} defines a ten-step performance evaluation
        methodology where step~6 (``Select the workload'') and step~7
        (``Design the experiments'') require workload-specific metrics
        and acceptance criteria, not uniform thresholds across
        heterogeneous workloads.\\
        \emph{Reference:} Jain, R., \emph{The Art of Computer Systems
        Performance Analysis: Techniques for Experimental Design,
        Measurement, Simulation, and Modeling}, Wiley, 1991.
        ISBN: 978-0-471-50336-1.

  \item \textbf{Gregg (2013, 2020)} describes \emph{active benchmarking}:
        analyzing system behavior during benchmark execution to identify
        the true limiter, rather than relying solely on pass/fail
        thresholds. The USE Method (Utilization, Saturation, Errors)
        provides a systematic approach to identifying bottlenecks---in
        this case, R's single-threaded CPU utilization reaching
        saturation under concurrent load. Gregg also warns against
        \emph{passive benchmarking} (running benchmarks without
        understanding what they measure) as an anti-method.\\
        \emph{Reference:} Gregg, B., \emph{Systems Performance:
        Enterprise and the Cloud}, 2nd ed., Prentice Hall, 2020.
        ISBN: 978-0-13-682015-4.\\
        \emph{Reference:} Gregg, B., ``Thinking Methodically about
        Performance'', \emph{Communications of the ACM}, Vol.~56,
        No.~2, Feb 2013.\\
        \emph{Reference:} Gregg, B., ``The USE Method'',
        \url{https://www.brendangregg.com/usemethod.html}.

  \item \textbf{Grafana k6 documentation} supports per-tag
        thresholds using the syntax
        \texttt{metric\{tag\_name:tag\_value\}}, enabling per-app
        SLAs within a single test run. The k6 testing guide recommends
        starting with smoke tests (low load) before progressing to
        average-load and stress tests---the baseline-then-load approach
        implements this progression.\\
        \emph{Reference:} Grafana Labs, ``Thresholds'', k6 Documentation,
        \url{https://grafana.com/docs/k6/latest/using-k6/thresholds/}.\\
        \emph{Reference:} Grafana Labs, ``Average-load testing'', k6
        Documentation,
        \url{https://grafana.com/docs/k6/latest/testing-guides/test-types/load-testing/}.

  \item \textbf{The RED Method} (Wilkie, 2015) measures Rate, Errors,
        and Duration per service/microservice. The per-app thresholds
        follow this principle: each Shiny app is a distinct microservice
        with its own expected duration profile.\\
        \emph{Reference:} Wilkie, T., ``The RED Method: key metrics
        for microservices architecture'', 2015.
        \url{https://www.weave.works/blog/the-red-method-key-metrics-for-microservices-architecture/}.
\end{itemize}

% ------------------------------------------------------------------
\section{REST Architecture Benchmark Results}
% ------------------------------------------------------------------

All benchmarks were run on the Hetzner CX32 cluster (8 vCPU, 16\,GB RAM)
with 3 backend replicas per Shiny app. Thresholds are per-app,
calibrated at $10\times$ baseline for relay/propagation/lifecycle
operations and $50\times$ baseline for throughput (100 VUs).

\subsection{Baseline (Benchmark 00)}

Single-user, no concurrency, 5 repetitions per app.

\begin{center}
\begin{tabular}{lrrrrr}
\toprule
\textbf{App} & \textbf{POST} & \textbf{Relay} & \textbf{Save} & \textbf{Restore} & \textbf{Total} \\
\midrule
Calculator     &   12\,ms &  3\,ms &  7\,ms &  4\,ms &    73\,ms \\
Analytics      &   13\,ms &  3\,ms &  6\,ms &  6\,ms &    83\,ms \\
Advanced       &   10\,ms &  3\,ms &  8\,ms &  5\,ms &    73\,ms \\
DataExchange   &   11\,ms &  3\,ms &  8\,ms &  6\,ms &    87\,ms \\
MonteCarlo     &   51\,ms &  4\,ms &  8\,ms &  5\,ms &   117\,ms \\
Geospatial     &   11\,ms &  3\,ms &  6\,ms &  4\,ms &    73\,ms \\
ClimateAnomaly &   29\,ms &  4\,ms &  8\,ms &  6\,ms &   105\,ms \\
MLTrainer      & 14,239\,ms &  4\,ms &  8\,ms &  5\,ms & 14,300\,ms \\
\bottomrule
\end{tabular}
\end{center}

\noindent Six of eight apps complete the full relay cycle in under
90\,ms. MonteCarlo (117\,ms) is slightly slower due to its simulation
loop. MLTrainer (14.3\,s) trains a 500-tree RandomForest on 2000 rows
per request---this is computation cost, not architecture overhead.

\subsection{Benchmark 01 — State Relay Latency (20 VUs)}

Measures the round-trip time of POST state $\to$ poll GET until the
processed result appears in Redis.

\begin{center}
\begin{tabular}{lrrrrl}
\toprule
\textbf{App} & \textbf{avg} & \textbf{med} & \textbf{p95} & \textbf{Threshold} & \textbf{Result} \\
\midrule
Calculator     &  2,795\,ms &    444\,ms & 13,681\,ms &    730\,ms & \textbf{FAIL} \\
Analytics      &  3,311\,ms &    598\,ms & 15,681\,ms &    830\,ms & \textbf{FAIL} \\
Advanced       &  5,157\,ms &  3,070\,ms & 13,478\,ms &    730\,ms & \textbf{FAIL} \\
DataExchange   &  4,352\,ms &  1,398\,ms & 14,128\,ms &    870\,ms & \textbf{FAIL} \\
MonteCarlo     &  2,552\,ms &  1,257\,ms &  8,013\,ms &  1,170\,ms & \textbf{FAIL} \\
Geospatial     &  3,770\,ms &    970\,ms & 12,705\,ms &    730\,ms & \textbf{FAIL} \\
ClimateAnomaly &  3,619\,ms &  1,082\,ms & 13,033\,ms &  1,050\,ms & \textbf{FAIL} \\
MLTrainer      & 34,917\,ms & 47,866\,ms & 56,670\,ms & 143,000\,ms & \textbf{PASS} \\
\bottomrule
\end{tabular}
\end{center}

\noindent All lightweight apps exceed their thresholds by 10--19$\times$.
The median latency (444\,ms--3,070\,ms) is acceptable for interactive
use, but the p95 tail (8--16\,s) is not. This tail is caused by
polling: under 20 VUs, each poll cycle (GET request) queues behind
other users' polls, compounding latency. MLTrainer passes because its
143\,s threshold absorbs the queuing overhead.

\subsection{Benchmark 02 — Collaboration Latency (30 VUs)}

Measures how long it takes for user B to see user A's state change
in a shared session.

\begin{center}
\begin{tabular}{lrrrrl}
\toprule
\textbf{App} & \textbf{avg} & \textbf{med} & \textbf{p95} & \textbf{Threshold} & \textbf{Result} \\
\midrule
Calculator     &  1,005\,ms &    123\,ms &  6,512\,ms &    730\,ms & \textbf{FAIL} \\
Analytics      &    401\,ms &    116\,ms &  1,498\,ms &    830\,ms & \textbf{FAIL} \\
Advanced       &  1,486\,ms &    176\,ms &  7,041\,ms &    730\,ms & \textbf{FAIL} \\
DataExchange   &  1,160\,ms &    198\,ms &  7,174\,ms &    870\,ms & \textbf{FAIL} \\
MonteCarlo     &  1,684\,ms &    217\,ms &  9,027\,ms &  1,170\,ms & \textbf{FAIL} \\
Geospatial     &    853\,ms &    112\,ms &  6,113\,ms &    730\,ms & \textbf{FAIL} \\
ClimateAnomaly &    820\,ms &     95\,ms &  5,306\,ms &  1,050\,ms & \textbf{FAIL} \\
MLTrainer      &  1,528\,ms &    647\,ms &  7,750\,ms & 143,000\,ms & \textbf{PASS} \\
\bottomrule
\end{tabular}
\end{center}

\noindent Median propagation is fast (95--647\,ms) but p95 reaches
6--9\,s. The same polling overhead applies: user B polls repeatedly
until user A's update appears, and each poll queues behind other
users' requests.

\subsection{Benchmark 03 — Save \& Restore Latency (15 VUs)}

\begin{center}
\begin{tabular}{lrrrrl}
\toprule
\textbf{App} & \textbf{Save avg} & \textbf{Save p95} & \textbf{Restore avg} & \textbf{Restore p95} & \textbf{Result} \\
\midrule
Calculator     & 1,673\,ms &  6,277\,ms & 2,194\,ms & 10,503\,ms & \textbf{FAIL} \\
Analytics      & 2,732\,ms &  7,070\,ms & 1,666\,ms &  9,838\,ms & \textbf{FAIL} \\
Advanced       & 3,038\,ms & 10,999\,ms & 2,874\,ms & 13,383\,ms & \textbf{FAIL} \\
DataExchange   & 3,433\,ms & 12,797\,ms & 1,375\,ms &  5,965\,ms & \textbf{FAIL} \\
MonteCarlo     & 3,064\,ms & 10,642\,ms & 2,858\,ms & 11,093\,ms & \textbf{FAIL} \\
Geospatial     & 3,346\,ms & 10,005\,ms & 1,667\,ms &  7,095\,ms & \textbf{FAIL} \\
ClimateAnomaly & 3,215\,ms & 10,160\,ms & 1,100\,ms &  5,235\,ms & \textbf{FAIL} \\
MLTrainer      & 3,927\,ms & 15,080\,ms & 2,566\,ms &  9,481\,ms & \textbf{FAIL} \\
\bottomrule
\end{tabular}
\end{center}

\noindent Save and restore operations go to PostgreSQL, which handles
concurrency better than R's single-threaded Plumber. However, the
save operation first triggers a state relay through the R backend
(to ensure the latest state is persisted), which introduces the same
queuing overhead. Even MLTrainer fails here because the save threshold
(80\,ms) is based on the database operation alone, not the full relay.

\subsection{Benchmark 04 — Throughput (100 VUs)}

\begin{center}
\begin{tabular}{lrrrrl}
\toprule
\textbf{App} & \textbf{avg} & \textbf{med} & \textbf{p95} & \textbf{Threshold} & \textbf{Result} \\
\midrule
Calculator     & 15,115\,ms & 13,845\,ms & 32,773\,ms &    600\,ms & \textbf{FAIL} \\
Analytics      & 16,869\,ms & 14,178\,ms & 36,021\,ms &    650\,ms & \textbf{FAIL} \\
Advanced       & 15,531\,ms & 13,753\,ms & 33,801\,ms &    500\,ms & \textbf{FAIL} \\
DataExchange   & 16,475\,ms & 13,429\,ms & 36,438\,ms &    550\,ms & \textbf{FAIL} \\
MonteCarlo     & 14,737\,ms & 12,882\,ms & 35,069\,ms &  2,550\,ms & \textbf{FAIL} \\
Geospatial     & 15,515\,ms & 13,747\,ms & 31,577\,ms &    550\,ms & \textbf{FAIL} \\
ClimateAnomaly & 16,296\,ms & 13,736\,ms & 36,845\,ms &  1,450\,ms & \textbf{FAIL} \\
MLTrainer      & 18,606\,ms & 16,697\,ms & 40,553\,ms & 711,950\,ms & \textbf{PASS} \\
\bottomrule
\end{tabular}
\end{center}

\noindent At 100 VUs with 3 replicas ($\sim$33 users per replica),
all lightweight apps average 13--17\,s per request. This is
$\sim$1000\times$ the baseline, far exceeding the $50\times$ threshold.
The R backend is fully saturated: each replica processes requests
sequentially, and the queue depth reaches 30+ requests. MLTrainer
passes trivially because its 712\,s threshold is unreachable.

\subsection{Benchmark 05 — Data Loss}

\textbf{Result: PASS.} Zero data loss across all apps. 3,169 states
sent, 3,150 received (19 lost to HTTP transport errors, not data
corruption). No state was silently dropped or overwritten.

\subsection{Benchmark 06 — Cross-Contamination}

\textbf{Result: PASS.} Zero contamination events. 108 isolation checks
passed. No state from one session leaked into another.

\subsection{Benchmark 07 — Session Lifecycle (50 VUs)}

\begin{center}
\begin{tabular}{lrrrrl}
\toprule
\textbf{App} & \textbf{avg} & \textbf{med} & \textbf{p95} & \textbf{Threshold} & \textbf{Result} \\
\midrule
Calculator     & 5,303\,ms & 3,208\,ms & 17,649\,ms &    730\,ms & \textbf{FAIL} \\
Analytics      & 5,122\,ms & 3,044\,ms & 16,823\,ms &    830\,ms & \textbf{FAIL} \\
Advanced       & 4,672\,ms & 2,791\,ms & 13,865\,ms &    730\,ms & \textbf{FAIL} \\
DataExchange   & 4,769\,ms & 2,844\,ms & 16,846\,ms &    870\,ms & \textbf{FAIL} \\
MonteCarlo     & 5,316\,ms & 3,376\,ms & 14,458\,ms &  1,170\,ms & \textbf{FAIL} \\
Geospatial     & 5,096\,ms & 3,125\,ms & 17,209\,ms &    730\,ms & \textbf{FAIL} \\
ClimateAnomaly & 5,080\,ms & 3,136\,ms & 18,697\,ms &  1,050\,ms & \textbf{FAIL} \\
MLTrainer      & 7,958\,ms & 5,309\,ms & 25,797\,ms & 143,000\,ms & \textbf{PASS} \\
\bottomrule
\end{tabular}
\end{center}

\noindent Lifecycle operations (create session, invite, join, post
state, save) each go through the Spring Boot backend. At 50 VUs,
the combined load from all operations saturates the backend.

\subsection{Benchmark 08 — Multi-User Collaboration (3 VUs)}

\begin{center}
\begin{tabular}{lrrrrl}
\toprule
\textbf{App} & \textbf{avg} & \textbf{med} & \textbf{p95} & \textbf{Threshold} & \textbf{Result} \\
\midrule
Calculator     &   24\,ms &   13\,ms &   49\,ms &    730\,ms & \textbf{PASS} \\
Analytics      &   18\,ms &   12\,ms &   32\,ms &    830\,ms & \textbf{PASS} \\
Advanced       &   80\,ms &   86\,ms &   87\,ms &    730\,ms & \textbf{PASS} \\
DataExchange   &   20\,ms &   15\,ms &   33\,ms &    870\,ms & \textbf{PASS} \\
MonteCarlo     &   16\,ms &   14\,ms &   26\,ms &  1,170\,ms & \textbf{PASS} \\
Geospatial     &   58\,ms &   54\,ms &   83\,ms &    730\,ms & \textbf{PASS} \\
ClimateAnomaly &   17\,ms &   21\,ms &   21\,ms &  1,050\,ms & \textbf{PASS} \\
MLTrainer      &   22\,ms &   16\,ms &   36\,ms & 143,000\,ms & \textbf{PASS} \\
\bottomrule
\end{tabular}
\end{center}

\noindent With only 3 VUs, all apps pass comfortably. Propagation
latency is 12--86\,ms, well within thresholds. This confirms that
the architecture works correctly at low concurrency---the failures
in other benchmarks are load-induced, not architectural.

\subsection{Benchmark 09 — WebSocket Presence (50 VUs)}

\textbf{Result: PASS.} WebSocket connect avg = 4\,ms (p95 = 5\,ms).
JOIN roundtrip avg = 6\,ms (p95 = 8\,ms). The STOMP-over-WebSocket
presence system handles 50 concurrent users with sub-10\,ms latency,
demonstrating that push-based communication scales where polling
does not.

\subsection{Summary}

\begin{center}
\begin{tabular}{llll}
\toprule
\textbf{\#} & \textbf{Benchmark} & \textbf{Result} & \textbf{Key Finding} \\
\midrule
00 & Baseline           & ---  & 6 apps $<$90\,ms, MLTrainer 14.3\,s \\
01 & State Relay        & 1/8 pass & Polling tail latency 8--16\,s at 20 VUs \\
02 & Collaboration      & 1/8 pass & Median fast (95--647\,ms), p95 slow (5--9\,s) \\
03 & Save/Restore       & 0/8 pass & Save triggers relay, inherits queuing \\
04 & Throughput         & 1/8 pass & 13--17\,s avg at 100 VUs ($\sim$1000$\times$ baseline) \\
05 & Data Loss          & \textbf{PASS} & 0\% loss \\
06 & Cross-Contamination & \textbf{PASS} & 0 events \\
07 & Session Lifecycle  & 1/8 pass & 13--18\,s p95 at 50 VUs \\
08 & Multi-User         & \textbf{8/8 pass} & 12--86\,ms at 3 VUs \\
09 & WebSocket          & \textbf{PASS} & 4--6\,ms avg, push scales \\
\bottomrule
\end{tabular}
\end{center}

\noindent \textbf{Data integrity is perfect}: zero data loss, zero
cross-contamination. The architecture is functionally correct.

\noindent \textbf{Latency degrades under load}: all polling-based
benchmarks fail their per-app thresholds at 20+ VUs. The bottleneck
is R's single-threaded request processing combined with the polling
overhead of REST state synchronization.

\noindent \textbf{Push-based communication scales}: WebSocket presence
(benchmark 09) handles 50 VUs at 6\,ms, and multi-user collaboration
(benchmark 08) passes at 3 VUs. The architecture works at low
concurrency; the degradation is load-induced.

\noindent \textbf{Compute-heavy apps are unaffected}: MLTrainer passes
every benchmark because its 14\,s computation cost dwarfs the
architecture overhead. For scientific workloads involving long-running
computations, the choice of messaging architecture is irrelevant.

% ------------------------------------------------------------------
\section{Climate Anomaly Checkpoint Enrichment}
% ------------------------------------------------------------------

The Climate Anomaly app uses a shared Kubernetes volume (PersistentVolumeClaim)
for data exchange: the R backend processes an uploaded CSV, writes the result to
\texttt{/app/shared\_data/processed\_summary.csv}, and sends a Kafka (or REST)
message containing only a file reference (\texttt{"file": "processed\_summary.csv"}).
The Shiny frontend reads the file from disk when it receives the message.

This design is intentional --- it demonstrates the shared-volume communication
pattern, which is architecturally distinct from the Data Exchange app's
inline-JSON pattern (where the dataset is embedded directly in the message).
Both patterns are benchmarked as part of the thesis comparison.

\subsection{The Overwrite Problem}

The file \texttt{processed\_summary.csv} is a singleton: every upload overwrites
it. When a checkpoint is saved, the \texttt{SessionStateMonitor} (Kafka) or
\texttt{SessionStateMonitor} (REST/Redis) captures the latest message, which
contains only the file reference. If another user uploads new data before the
checkpoint is restored, the file on disk no longer matches the checkpoint's
original state. Cross-user restore (RQ5) fails because the Shiny app reads
stale or unrelated data.

\subsection{Solution: Save-Time Enrichment, Restore-Time File Reconstruction}

The fix preserves the shared-volume pattern during normal operation while making
checkpoints self-contained:

\begin{enumerate}[nosep]
  \item \textbf{Save time:} Spring Boot's \texttt{StateController} (Kafka) or
        \texttt{CollaborationController} (REST) detects \texttt{CLIMATE\_READY}
        messages, reads the referenced CSV from the shared volume (mounted via
        PVC at \texttt{/app/shared\_data}), parses it into a list of row objects,
        and injects it as a \texttt{"dataset"} field in the checkpoint JSON
        before persisting to the database.
  \item \textbf{Restore time:} Before publishing the restored state, Spring Boot
        extracts the \texttt{"dataset"} field, reconstructs the CSV file, writes
        it back to the shared volume at the original path, and removes
        \texttt{"dataset"} from the message payload. The Shiny frontend receives
        the same file-reference message it expects and reads from disk as usual.
\end{enumerate}

This approach required mounting the shared-data PVC on the Spring Boot container
(previously only the Shiny containers had access). On a single-node k3s cluster
with \texttt{ReadWriteOnce} access mode, multiple pods on the same node can
mount the same PVC without issue.

\subsection{Why Not Inline Data in Normal Operation?}

Sending the dataset inline during normal operation (as Data Exchange does) would
work but would eliminate the architectural distinction between the two apps.
The Climate Anomaly app exists specifically to benchmark the shared-volume
pattern. The enrichment-at-save approach preserves this distinction: normal
message flow uses file references (benchmarkable), while checkpoints embed the
data (reproducible).

\subsection{Known Limitation: CSV Parser}

The Java CSV parser in \texttt{enrichClimateCheckpoint()} and
\texttt{restoreClimateFile()} uses \texttt{String.split(",")} to tokenize rows.
This does not handle RFC~4180 quoted fields correctly: a value like
\texttt{"New York, NY"} would be split into two fields instead of one.

For the Climate Anomaly data this is safe. The columns produced by the R backend
are \texttt{SiteID} (e.g., \texttt{SITE\_A}), \texttt{Date}
(\texttt{2026-06-01}), two numeric means, and \texttt{Heatwave\_Anomaly}
(\texttt{YES}/\texttt{NO}) --- none of which contain commas. However, if the
schema were extended to include free-text fields, the parser would silently
corrupt data.

A production implementation should use a proper CSV library such as OpenCSV
(\texttt{com.opencsv.CSVReader}) which handles quoting, escaping, and
multi-line values correctly. For the thesis scope, the simple parser is
sufficient and avoids adding a new dependency.

\subsection{Known Limitation: Code Duplication}

The \texttt{enrichClimateCheckpoint()} and \texttt{restoreClimateFile()} methods
are duplicated across three Java files:

\begin{enumerate}[nosep]
  \item \texttt{Kafka/StateController.java} --- solo-mode save/restore
  \item \texttt{REST/StateController.java} --- solo-mode save/restore
  \item \texttt{REST/CollaborationController.java} --- session-mode save/restore
\end{enumerate}

This violates the DRY (Don't Repeat Yourself) principle. A cleaner design would
extract these methods into a shared Spring \texttt{@Component} (e.g.,
\texttt{ClimateCheckpointHelper}) injected into both controllers. This was not
done because the methods are small (approximately 40 lines each), the risk of
divergence is low in a thesis project with a single developer, and the
refactoring would add indirection without functional benefit. If the checkpoint
enrichment logic were extended to additional apps, extracting a shared component
would become necessary.

\subsection{Known Limitation: Singleton File and Concurrent Restores}

The restore operation writes the embedded data back to the shared volume at the
original filename (\texttt{processed\_summary.csv}). If two users restore
different Climate Anomaly checkpoints simultaneously, the second write
overwrites the first, and one user may see incorrect data. This is the same
singleton-file problem that motivated the enrichment in the first place, now
occurring at restore time instead of upload time.

A production fix would generate a unique filename per restore (e.g.,
\texttt{processed\_summary\_\{stateId\}.csv}) and update the \texttt{file}
field in the payload accordingly. For the thesis, concurrent restores are
unlikely in the test scenarios and benchmarks.

\subsection{Future Work: Object Storage (S3) Instead of PVC}

The current shared-volume approach uses a Kubernetes PersistentVolumeClaim (PVC)
with \texttt{ReadWriteOnce} access mode. This works on a single-node cluster
but has fundamental limitations:

\begin{itemize}[nosep]
  \item \textbf{Single-node constraint:} \texttt{ReadWriteOnce} PVCs can only
        be mounted by pods on the same node. Scaling to multiple nodes requires
        \texttt{ReadWriteMany}, which needs a network filesystem (NFS, CephFS,
        or a cloud provider's shared storage).
  \item \textbf{No versioning:} Files are overwritten in place. The enrichment
        mechanism works around this by embedding data in checkpoints, but the
        underlying problem remains.
  \item \textbf{No access control:} Any pod with the volume mount can read or
        write any file. There is no per-user or per-session isolation.
  \item \textbf{Tight coupling:} Spring Boot needed a PVC mount added to its
        deployment just to read CSV files at checkpoint save time. Every new
        consumer of shared data requires a deployment change.
\end{itemize}

An object storage service (e.g., Amazon S3, MinIO, or Google Cloud Storage)
would address all of these:

\begin{itemize}[nosep]
  \item \textbf{Multi-node:} S3 is network-accessible; no volume mounts needed.
        Any pod can read/write via HTTP, regardless of node placement.
  \item \textbf{Versioning:} S3 supports object versioning natively. Each upload
        creates a new version; checkpoints can reference a specific version ID
        instead of embedding the data.
  \item \textbf{Access control:} IAM policies and pre-signed URLs provide
        fine-grained access control per object, per user, or per session.
  \item \textbf{Loose coupling:} Services access data via an API client (e.g.,
        the AWS SDK), not filesystem mounts. Adding a new consumer requires no
        deployment changes.
  \item \textbf{No enrichment needed:} With versioned objects, the checkpoint
        stores the S3 key and version ID. The file is immutable at that version,
        so restore simply reads the correct version --- no embedding, no file
        reconstruction, no race conditions.
\end{itemize}

The trade-off is operational complexity: S3 requires either a cloud account or
a self-hosted MinIO instance, plus credential management. For a thesis
deployment on a single Hetzner node, the PVC approach is simpler and sufficient.
For a production NaaVRE integration, S3 or MinIO would be the recommended
approach.

\subsection{Future Work: Kafka-Level Compression}

For larger datasets, Kafka's built-in producer compression
(\texttt{compression.type = gzip}) could reduce message sizes transparently
without any application code changes. Similarly, a hybrid approach with a size
threshold (inline for small datasets, file reference for large ones) could be
implemented, but was not needed for the Climate Anomaly summary which is
typically dozens of rows after aggregation.

% ------------------------------------------------------------------
\section{Future Work: A Background-Thread Kafka Library for R}
% ------------------------------------------------------------------

A purpose-built R Kafka library could significantly reduce consumer
latency, though not eliminate the single-thread constraint. The design
would be:

\begin{enumerate}
  \item \textbf{Core in Rust or C++}, linked against librdkafka.
        Rust's \texttt{rdkafka} crate or librdkafka's C API directly.
  \item \textbf{Background consumer thread} that runs
        \texttt{rd\_kafka\_consumer\_poll()} in a loop, independent of R's
        event loop. Messages are pushed into a lock-free, thread-safe
        ring buffer.
  \item \textbf{Non-blocking R interface}: \texttt{consumer\$poll()} returns
        immediately with whatever messages are buffered, or
        \texttt{NULL} if none. No timeout, no blocking.
  \item \textbf{Callback registration}: optionally, the library invokes
        an R function when a message matching a filter (topic, key,
        partition) arrives, scheduled via \texttt{later::later()} to
        integrate with Shiny's event loop.
  \item \textbf{Partition-aware consumption}: expose
        \texttt{consumer\$assign(topic, partition)} so applications can
        subscribe to specific partitions rather than scanning an entire
        topic.
\end{enumerate}

The R-side integration with Shiny would replace \texttt{reactivePoll()}
with a \texttt{later}-based callback:

\begin{lstlisting}[language=R]
# Hypothetical API
consumer <- KafkaConsumer$new(brokers = "kafka:9092", group = "shiny")
consumer$subscribe("output", key_filter = session_id)
consumer$on_message(function(msg) {
  data <- jsonlite::fromJSON(msg$value)
  shared_state(data)  # update reactive value
})
\end{lstlisting}

This eliminates the poll-wait-discard cycle entirely. The background
thread handles all Kafka I/O; R only executes when a relevant message
arrives.

\subsection{Feasibility}

The \texttt{extendr} package (Rust $\to$ R bindings) and the
\texttt{cpp11} package (C++ $\to$ R bindings) provide mature foreign
function interfaces. The \texttt{httpuv} package already proves that
background-thread I/O with R callbacks is viable in production
(it powers every Shiny application).

The main engineering challenge is thread safety: R's C API is not
thread-safe, so the background thread cannot call R functions directly.
All R interaction must be marshalled through \texttt{later::later()},
which schedules execution on R's main thread at the next idle point.
This is a solved problem---\texttt{httpuv} and \texttt{promises} both
use this pattern.

Estimated effort: 2--3 months for a working prototype, assuming
familiarity with Rust/C++ and librdkafka.

% ------------------------------------------------------------------
\section{Future Work: Reusable R Package for State Sharing}
% ------------------------------------------------------------------

A natural extension of ShinySwarm's checkpoint mechanism is a standalone R package
(e.g., \texttt{shinyswarm}) that any Shiny application can use for full-state
serialisation and restoration, without requiring the full microservice architecture
(Spring Boot, Redis, Kafka).

\subsection{Motivation}

Every Shiny app in ShinySwarm follows the same pattern: capture state (inputs
plus computed reactive values), serialise it to JSON, and later restore it by
pushing values back into the app. Currently this logic is implemented
individually in each of the eight apps, with each one manually building JSON
payloads with app-specific field names. A library would standardise this into
two function calls.

Existing Shiny state-sharing mechanisms (\texttt{shinyURL}, bookmarking)
transfer only inputs, forcing recomputation of all outputs. For non-deterministic
analyses (e.g., Monte Carlo simulations without \texttt{set.seed()}), recomputing
produces different results. A full-state package preserves outputs alongside
inputs, guaranteeing exact reproducibility.

\subsection{API Design}

The package exposes two core functions. The key design decision is that the
library is \textbf{transport-agnostic}: it produces and consumes JSON strings.
How the JSON is transmitted or persisted (REST API, Kafka, Redis, file, database)
is the caller's responsibility.

\subsubsection{\texttt{capture\_state()} --- Serialise Everything to JSON}

\begin{verbatim}
capture_state <- function(input, ...) {
  state <- list()

  # 1. Capture all input values
  input_names <- names(input)
  state$inputs <- lapply(
    setNames(input_names, input_names),
    function(name) input[[name]]
  )

  # 2. Capture registered reactive values
  #    Passed explicitly because R cannot discover
  #    reactiveVal objects automatically
  extras <- list(...)
  state$reactives <- lapply(extras, function(rv) {
    if (is.reactivevalues(rv)) {
      reactiveValuesToList(rv)
    } else {
      rv()  # reactiveVal -- call to get value
    }
  })

  state$timestamp <- as.numeric(Sys.time())
  toJSON(state, auto_unbox = TRUE, force = TRUE)
}
\end{verbatim}

The developer passes reactive values by name: \texttt{capture\_state(input,
result = result, shared\_df = shared\_df)}. The function reads all
\texttt{input} values automatically via \texttt{names(input)}, then reads each
registered reactive by calling it (for \texttt{reactiveVal}) or converting it
(for \texttt{reactiveValues}).

\subsubsection{\texttt{restore\_state()} --- Push Everything Back}

\begin{verbatim}
restore_state <- function(session, state_json, ...) {
  state <- fromJSON(state_json, simplifyVector = FALSE)

  # 1. Restore inputs
  for (name in names(state$inputs)) {
    value <- state$inputs[[name]]
    session$sendInputMessage(name, list(value = value))
  }

  # 2. Restore reactive values
  extras <- list(...)
  for (name in names(state$reactives)) {
    if (name %in% names(extras)) {
      rv <- extras[[name]]
      saved_value <- state$reactives[[name]]
      if (is.reactivevalues(rv)) {
        for (key in names(saved_value)) {
          rv[[key]] <- saved_value[[key]]
        }
      } else {
        rv(saved_value)  # reactiveVal -- set directly
      }
    }
  }
}
\end{verbatim}

Inputs are restored via \texttt{session\$sendInputMessage()}, which is the
low-level Shiny mechanism for pushing a value into any input widget. Reactive
values are set directly by calling the setter (for \texttt{reactiveVal}) or
assigning keys (for \texttt{reactiveValues}).

\subsection{Usage Example: Calculator App}

With the library, the Calculator app's save/restore logic reduces to:

\begin{verbatim}
library(shiny)
library(shinyswarm)

server <- function(input, output, session) {
  result <- reactiveVal(0)

  observeEvent(input$calculate, {
    result(input$num1 + input$num2)
  })

  output$result <- renderText({ result() })

  # Save: one line
  observeEvent(input$save_btn, {
    json <- capture_state(input, result = result)
    send_to_backend(json)  # REST, Kafka, file -- caller decides
  })

  # Restore: one line
  observeEvent(input$restore_btn, {
    json <- get_from_backend()
    restore_state(session, json, result = result)
  })
}
\end{verbatim}

Compare this to the current approach where each app manually constructs JSON
payloads with app-specific field names and handles each input individually.

\subsection{Transport Agnosticism}

The library deliberately excludes any transport or storage logic. The same
\texttt{capture\_state()} output can be sent via any mechanism:

\begin{verbatim}
# Via REST API
httr::POST("http://backend/state", body = json)

# Via Kafka
producer$produce("output", json, key = session_id)

# Via file
writeLines(json, "checkpoint.json")

# Via database
DBI::dbExecute(con,
  "INSERT INTO states (data) VALUES (?)", list(json))
\end{verbatim}

This makes the package useful to any Shiny developer, not just ShinySwarm users.
A researcher running a standalone Shiny app on their laptop could use it to save
and restore state to a local file, while a ShinySwarm deployment uses the same
functions with Kafka or REST transport.

\subsection{Implementation Challenges}

\begin{enumerate}[nosep]
  \item \textbf{No reactive discovery API.} R's reactive system does not expose
        a ``get all reactive values'' function. \texttt{reactiveVal} and
        \texttt{reactiveValues} objects must be explicitly registered by the app
        developer --- they cannot be discovered automatically. This is a
        fundamental limitation of Shiny's design, not something the library can
        work around.

  \item \textbf{\texttt{fileInput} cannot be restored.}
        \texttt{session\$sendInputMessage()} does not work for file inputs.
        The uploaded file existed at a temporary path that no longer exists
        after the session ends. A workaround would store the file content
        (base64-encoded) in the checkpoint JSON, then on restore write it to a
        new temporary file and inject the path. This is similar to the Climate
        Anomaly enrichment approach already implemented in ShinySwarm.

  \item \textbf{Rendered outputs are server-side objects.} Plots
        (\texttt{renderPlot}), tables (\texttt{renderTable}), and other outputs
        cannot be serialised directly. The library stores the \emph{data} that
        generates them (via the registered reactive values), not the rendered
        output itself. On restore, the reactive dependency chain triggers
        re-rendering automatically. This is correct for deterministic outputs
        but means the restore is not instantaneous --- the app must recompute
        the rendering.

  \item \textbf{Shiny modules use namespaced IDs.} A module input
        \texttt{"slider"} becomes \texttt{"module1-slider"} internally.
        \texttt{capture\_state()} sees the namespaced ID, but
        \texttt{restore\_state()} must route it through the module's session
        object, not the parent session. Supporting modules would require either
        a recursive capture/restore across session hierarchies, or requiring
        each module to call \texttt{capture\_state()} independently.

  \item \textbf{Non-deterministic outputs.} Monte Carlo simulations without
        \texttt{set.seed()} produce different results each run. If only inputs
        are stored (as \texttt{shinyURL} and bookmarking do), restoring and
        recomputing gives different outputs. This is the primary motivation for
        storing reactive values (computed results) alongside inputs --- the
        checkpoint preserves the exact output, not just the recipe to
        regenerate it.
\end{enumerate}

\subsection{Typed State Wrappers}

The basic \texttt{capture\_state()} API handles scalar inputs and reactive values,
but Shiny apps contain different \emph{types} of state that require different
capture strategies:

\begin{center}
\begin{tabular}{lll}
\textbf{State type} & \textbf{Example} & \textbf{Capture strategy} \\
\hline
Scalar inputs & \texttt{input\$num1 = 42} & \texttt{input[[name]]} \\
Reactive dataframes & \texttt{shared\_df()} & \texttt{rv()} $\to$ JSON array \\
File references & \texttt{processed\_summary.csv} & Read file $\to$ base64 encode \\
Binary R objects & Trained ML model & \texttt{serialize()} $\to$ base64 \\
\end{tabular}
\end{center}

Rather than creating separate functions (\texttt{capture\_state\_file()},
\texttt{capture\_state\_json()}, etc.), the library uses \textbf{type-hinted
wrappers} that tell \texttt{capture\_state()} how to handle each piece of state.
The output is always a single JSON string --- the wrappers control what goes
\emph{inside} that JSON.

\subsubsection{Wrapper Functions}

\begin{verbatim}
# Wrapper: capture a file from disk (read + base64)
state_file <- function(path) {
  structure(list(path = path), class = "shinyswarm_file")
}

# Wrapper: capture an arbitrary R object via serialize()
state_rds <- function(obj) {
  structure(list(obj = obj), class = "shinyswarm_rds")
}
\end{verbatim}

The developer registers each reactive with the appropriate wrapper:

\begin{verbatim}
json <- capture_state(
  input,
  result    = result,           # reactiveVal -- just call rv()
  shared_df = shared_df,        # dataframe -- toJSON()
  csv_file  = state_file(       # file -- read + base64
    "/app/shared_data/processed_summary.csv"
  ),
  model     = state_rds(        # R object -- serialize() + base64
    trained_model
  )
)
\end{verbatim}

\subsubsection{Capture Dispatch}

Inside \texttt{capture\_state()}, the function dispatches based on the class
of each registered value:

\begin{verbatim}
for (name in names(extras)) {
  rv <- extras[[name]]

  if (inherits(rv, "shinyswarm_file")) {
    raw <- readBin(rv$path, "raw", file.info(rv$path)$size)
    state$reactives[[name]] <- list(
      type = "file",
      filename = basename(rv$path),
      content = base64enc::base64encode(raw)
    )

  } else if (inherits(rv, "shinyswarm_rds")) {
    state$reactives[[name]] <- list(
      type = "rds",
      content = base64enc::base64encode(
        serialize(rv$obj, NULL)
      )
    )

  } else if (is.reactivevalues(rv)) {
    state$reactives[[name]] <- list(
      type = "reactiveValues",
      content = reactiveValuesToList(rv)
    )

  } else {
    state$reactives[[name]] <- list(
      type = "reactiveVal",
      content = rv()
    )
  }
}
\end{verbatim}

\subsubsection{Restore Dispatch}

\texttt{restore\_state()} reads the \texttt{type} field and reverses the
operation:

\begin{verbatim}
for (name in names(state$reactives)) {
  saved <- state$reactives[[name]]
  rv <- extras[[name]]

  if (saved$type == "file") {
    raw <- base64enc::base64decode(saved$content)
    path <- file.path("/app/shared_data", saved$filename)
    writeBin(raw, path)

  } else if (saved$type == "rds") {
    obj <- unserialize(
      base64enc::base64decode(saved$content)
    )
    rv(obj)

  } else if (saved$type == "reactiveValues") {
    for (key in names(saved$content)) {
      rv[[key]] <- saved$content[[key]]
    }

  } else {
    rv(saved$content)
  }
}
\end{verbatim}

\subsubsection{Relationship to Climate Anomaly Enrichment}

This typed-wrapper design generalises the Climate Anomaly checkpoint enrichment
already implemented in ShinySwarm. The current implementation performs the
file-read-and-embed logic in Java (Spring Boot's \texttt{enrichClimateCheckpoint()}
and \texttt{restoreClimateFile()} methods). With the R package, this logic would
move into the Shiny app itself via \texttt{state\_file()}, eliminating the need
for Spring Boot to mount the shared volume, parse CSV files, or reconstruct
them on restore. The entire enrichment mechanism would reduce to a single
wrapper call in the app's save handler.

\subsection{Package Structure}

\begin{verbatim}
shinyswarm/
  DESCRIPTION
  NAMESPACE
  R/
    capture_state.R     # capture_state()
    restore_state.R     # restore_state()
    wrappers.R          # state_file(), state_rds()
    utils.R             # base64 helpers
  man/                  # auto-generated roxygen docs
  tests/
    testthat/
      test-capture.R
      test-restore.R
      test-wrappers.R
\end{verbatim}

The package could be published on CRAN or as a GitHub-installable package
(\texttt{remotes::install\_github("...")}), making ShinySwarm's reproducibility
contribution accessible to the broader R~Shiny community without requiring
Kubernetes or microservice infrastructure.

% ------------------------------------------------------------------
\section{Future Work: Integration with NaaVRE}
% ------------------------------------------------------------------

NaaVRE (Notebook-as-a-Virtual-Research-Environment) is a JupyterLab-based
platform developed at the University of Amsterdam that allows researchers
to containerize notebook cells, compose them into workflows, and execute
them on Kubernetes via Argo Workflows. Its architecture consists of
JupyterLab frontend extensions communicating with backend API services
through a central communication hub, with Keycloak providing OIDC
authentication.

A natural extension of this thesis is embedding collaborative R Shiny
visualizations into NaaVRE as interactive workflow outputs. The ShinySwarm
architecture developed in this thesis already solves the core technical
challenges required for this integration.

\subsection{How NaaVRE Containers Communicate}

NaaVRE follows a microservice pattern on Kubernetes:

\begin{itemize}
  \item \textbf{UI components} are JupyterLab frontend extensions
        (TypeScript), communicating with backend services through a
        server extension called the NaaVRE Communicator.
  \item \textbf{API services} are independent FastAPI or Django
        microservices. They do not communicate with each other directly.
  \item \textbf{Workflow execution} uses Argo Workflows: each notebook
        cell is containerized into a Docker image and orchestrated as a
        DAG. Cells pass data via workflow parameters or shared storage.
  \item \textbf{Authentication} is centralized through Keycloak using
        OIDC/JWT tokens, validated by each API service independently.
\end{itemize}

The key distinction from ShinySwarm: NaaVRE's inter-container
communication is batch/pipeline-style (cell A produces output, cell B
consumes it), not real-time. There is no persistent WebSocket or
streaming channel between containers.

\subsection{Proposed Integration Architecture}

Collaborative Shiny applications can be integrated into NaaVRE as
interactive visualization services. Three approaches are possible:

\begin{enumerate}
  \item \textbf{Shiny as a NaaVRE UI component (recommended).}
        Deploy Shiny frontends as K8s services (as in this thesis) and
        embed them in JupyterLab via a custom extension using iframes.
        The JupyterLab extension replaces Angular as the host application,
        communicating with Shiny via \texttt{postMessage} for RBAC and
        STOMP-over-WebSocket for presence---the same mechanisms already
        implemented and tested in ShinySwarm.

  \item \textbf{Shiny as a workflow output viewer.}
        An Argo workflow cell produces data and writes it to shared
        storage (e.g.\ a PersistentVolumeClaim or S3 bucket). A Shiny
        app reads from that storage to render interactive visualizations.
        No real-time synchronization is needed; the Shiny app simply
        displays the latest output.

  \item \textbf{Shiny inside JupyterLab via \texttt{jupyter-shiny-proxy}.}
        The existing \texttt{jupyter-shiny-proxy} package runs Shiny
        apps as JupyterLab server extensions. Each user gets a dedicated
        R process. However, this bypasses NaaVRE's containerization model
        and does not support multi-user collaboration.
\end{enumerate}

\subsection{What ShinySwarm Already Provides}

The architecture developed in this thesis provides a ready-made
foundation for option~1:

\begin{itemize}
  \item \textbf{Iframe embedding with \texttt{postMessage} relay} for
        propagating role-based access control (RBAC) changes from the
        host application into the Shiny iframe. This mechanism is
        framework-agnostic---it works identically whether the host is
        Angular, React, or a JupyterLab extension.
  \item \textbf{STOMP-over-WebSocket presence} for real-time user
        awareness (who is viewing, who is editing). The Spring Boot
        backend handles this independently of the frontend framework.
  \item \textbf{State synchronization via REST API} using Redis as the
        state store and PostgreSQL for persistence. The Shiny frontends
        poll \texttt{/api/collab/\{sessionId\}/state} via HTTP.
  \item \textbf{Session lifecycle management} (create, join, leave, save,
        restore) exposed as REST endpoints, directly consumable by a
        JupyterLab extension.
\end{itemize}

The migration effort from Angular to JupyterLab is a frontend change,
not an architectural one. The Shiny services, Spring Boot backend, and
state synchronization layer remain unchanged.

\end{document}
